\clearpage
\fancyhead[L]{\small ZACH. TABLES D'ABERRATION.}
\hypertarget{native-551-lower}{}\label{native-551-lower}
\begin{center}\small Göttingische gelehrte Anzeigen. Stück 196. Seite 1956..1960. 1812 December 7.\\[3pt]\rule{.72\linewidth}{.3pt}\end{center}

\emph{Nouvelles tables d'aberration et de nutation pour quatorze cent quatre étoiles, avec une table générale d'aberration pour les planètes et les comètes, précédées d'une instruction qui renferme l'explication de l'usage de ces tables, suivies de plusieurs nouvelles tables destinées à faciliter les calculs astronomiques, par le Baron} \textsc{de Zach}. De l'imprimerie de Madame Mine et Comp. Marseille. 52 und 136 Seiten in Octav.

Seit der Entdeckung der Aberration und Nutation haben die Astronomen auf vielfache Weise die lästige und dem beobachtenden Astronomen täglich wiederkehrende Berechnung dieser kleinen Ungleichheiten zu erleichtern gesucht. Man hat mancherlei mehr oder weniger glücklich angeordnete allgemeine Tafeln dafür, man hat für einzelne Sterne specielle Tafeln, welche die Aberration und Nutation nach der jedesmaligen Stellung der Sonne und des Mondsknoten unmittelbar geben. Die Tafeln \SourceLeaf{552}der letztern Gattung von \textsc{Mezger}, \textsc{von Zach} und \textsc{Cagnoli} sind in den Händen aller practischen Astronomen, beschränken sich aber noch auf eine zu kleine Anzahl von Sternen, als dass man nicht am häufigsten zu den allgemeinen Tafeln zurück zu kehren genöthigt wäre. Hr. \textsc{von Zach} hat jetzt in der vorliegenden Schrift für eine beträchtlich grössere Zahl von Sternen zwar keine eigentliche Specialtafeln geliefert, als welche viel zu voluminös und kostbar ausfallen würden, aber doch die Rechnung noch etwas mehr zu erleichtern gesucht, als es selbst durch die am zweckmässigsten eingerichteten allgemeinen Tafeln möglich ist. Die Idee davon ist sehr einfach und in Beziehung auf Aberration auch schon früher von andern Astronomen benutzt. Die Aberration der geraden Aufsteigung eines jeden Sterns lässt sich leicht in die Form bringen $m\sin(\odot+A)$, wo $m$ die grösste positive Aberration des Sterns, $\odot$ die jedesmalige Länge der Sonne, und $A$ das Complement derjenigen bestimmten Sonnenlänge zu $360^\circ$ ausdrückt, bei welcher die Aberration verschwindet und aus dem negativen in den positiven Werth übergeht. Die Grössen $m$ und $A$ hängen von der mittlern Stellung des Sterns ab, und da es bei ihnen auf sehr grosse Genauigkeit nicht ankommt, kann man sich der einmal bestimmten Werthe eine ziemliche Reihe von Jahren hindurch unverändert bedienen. Ganz eben so verhält es sich mit der Aberration der Declination, der Nutation in gerader Aufsteigung, und der Nutation in der Declination, nur dass $m$ und $A$ natürlich für jedes dieser Elemente andere Werthe haben. Diese vier Werthe von $A$ und die Logarithmen von der dazu gehörenden $m$ liefert Hr. \textsc{von Zach}, jene auf Minuten, diese auf vier Decimalen für 1404 der vornehmsten Sterne aus allen in Europa sichtbaren Sternbildern nach ihren mittlern Positionen für 1800, auf 94 Seiten, wodurch sich also das ganze Geschäft auf das Aufschlagen von vier Sinus-Logarithmen und vier Zahlen aus den Logarithmen reducirt. Man hat so die logarithmischen Tafeln an vier Stellen nachzuschlagen, also einmal mehr, als bei dem Gebrauche der \emph{allgemeinen} Tafeln von \textsc{Gauss}, aber doch zusammen weniger zu schreiben, als bei diesen, und wird also die neuen \textsc{von Zach}'schen Tafeln für alle Sterne, die man in ihnen findet, den allgemeinen gern vorziehen. In der übrigens sehr lehrreichen Einleitung des Hrn. \textsc{von Zach} hat sich bei der theoretischen Entwickelung eine kleine Unrichtigkeit eingeschlichen, indem dem Hülfsmittel $A$ eine falsche Bedeutung beigelegt ist. --- Nach diesen Tafeln für 1404 Sterne folgen die 36 so genannten \textsc{Maskelyne}'schen Sterne noch einmal, zuerst ihre geraden Aufsteigungen nach \textsc{Maskelyne}, nebst $A$ und $\log m$ in Zeit für Aberration und Nutation, sodann von denselben Sternen die Declinationen nach \textsc{Piazzi}, nebst den beiden sich darauf beziehenden Werthen von $A$ und $\log m$, alles für 1802. Wir bemerken hier einige, an sich freilich unbedeutende, Differenzen mit der ersten Tafel, welche sich aus der blossen Veränderung der Epoche nicht erklären lassen, sondern Rechnungsfehler zu sein scheinen. So ist z. B. bei der Declination von $\alpha$ Tauri \[ \begin{array}{r@{}l|c|c} \multicolumn{2}{c|}{} & \text{erste Tafel} & \text{zweite Tafel}\\ \text{Aberration}\; & \left\{\begin{array}{l}\log m\\ A\end{array}\right. & \begin{array}{c}0\,5806\\7^{\mathrm s}\;23^\circ\;26'\end{array} & \begin{array}{c}0.5772\\7^{\mathrm s}\;23^\circ\;15'\end{array}\\ \text{Nutation}\; & \left\{\begin{array}{l}\log m\\ A\end{array}\right. & \begin{array}{c}0.9673\\3^{\mathrm s}\;18^\circ\;18'\end{array} & \begin{array}{c}0.9682\\3^{\mathrm s}\;18^\circ\;6'\end{array} \end{array} \]

Eine ähnliche Tafel folgt hierauf für die Sterne $\alpha$ und $\beta$ im kleinen Bär nach ihren Stellungen in den Jahren 1790, 1800, 1810 und 1820.

Obgleich die Grössen $m$ und $A$ eigentlich nur für das Jahr gelten, für welches sie berechnet sind, so kann man doch, ohne einen erheblichen Fehler zu begehen, sie für eine ziemlich lange Reihe von Jahren, vor und nach jener Epoche beibehalten. Der Verf. gibt ein paar besondere Tafeln, um zu überschlagen, wie viel gerade Aufsteigung und Abweichung sich ändern dürfen, ehe der Fehler, welcher daraus entspringt, auf 1 Secunde wächst. Diese Tafel dient zugleich, zu beurtheilen, in wie fern \SourceLeaf{553}man für einen Stern, welcher sich nicht unter den 1404 der grossen Tafel findet, die Constanten eines benachbarten Sterns aus derselben zu gebrauchen sich erlauben dürfe. Einige hierauf folgende Tafeln für Aberration, Nutation, Präcession und Verbesserung des Mittags aus correspondirenden Sonnenhöhen, welche der Verf. auch schon an andern Orten gegeben hatte, übergehen wir hier als bekannt mit Stillschweigen. --- Bequem ist die Tafel XX, um die Sonnenlänge als Aberrationsargument für den Tag eines Jahrs zu finden, wofür man keinen astronomischen Kalender zur Hand hat; sie gibt an, wie viel in jedem einzelnen Jahre von 1700 bis 1827 die Epoche der mittlern Sonnenlänge grösser ist, als im Jahre 1803. Für den angezeigten Zweck genau genug, kann man diese Zahlen zugleich als die Unterschiede der wahren Längen an gleichnamigen Tagen betrachten, und so leicht die Sonnenlängen für jedes jener Jahre aus denen irgend eines andern, wofür man eben einen astronomischen Kalender zur Hand hat, ableiten. Eben so liefert Hr. \textsc{von Zach} hier noch eine für die Länge des Mondsknoten beim Anfang aller Jahre, von 1700 bis 1827, um die Ephemeriden für irgend eines dieser Jahre auf ein anderes übertragen zu können. Zum Schluss folgt noch eine Tafel zur Verwandlung der Sternzeit in mittlere Sonnenzeit, noch etwas bequemer eingerichtet, als die bekannten, zu demselben Zweck berechneten, kleinen Tafeln des Verf., welche in den Händen aller Astronomen sind. Es ist dabei zugleich eine kleine Tafel für denjenigen Theil der Nutation beigefügt, welcher von der Sonnenlänge abhängt; dieser hätte allerdings sogleich mit der Tafel XXII vereinigt werden können, wodurch die Rechnung noch ein wenig verkürzt wäre. Vermuthlich hat der Verf. es absichtlich unterlassen, damit man in dieser Tafel die reinen mittlern Sonnen-Rectascensionen habe; denn aus demselben Grunde scheint die Hinzufügung einer negativen Constante unterblieben zu sein, die man sonst anwendet, um die Monds-Nutation immer positiv zu machen.

Man sieht aus dieser kurzen Inhaltsanzeige, dass der verdiente Verf. sich durch die Sammlung von Tafeln abermals neue Ansprüche auf den Dank aller practischen Astronomen erworben habe, denen jede Erleichterung ihrer mühsamen, täglich wiederkehrenden, Rechnungen sehr willkommen sein muss.
